\input{../tex-inputs/latex-korrekturansicht-vorspann}

\section[Berta Zuckerkandl an Arthur Schnitzler, {[}23. 6. 1911?{]}]{L03995 Berta Zuckerkandl an Arthur Schnitzler, {[}23. 6. 1911?{]}}
\nopagebreak\mylabel{L03995v}
\rehead{ }\normalsize\beginnumbering\briefempfaengerindex{Schnitzler, Arthur@\textsc{Schnitzler, Arthur}!zzzZuckerkandl, Berta@\emph{von Berta Zuckerkandl}!1911-06-231@{{[}23. 6. 1911?{]}}|(be}
\toendnotes[C]{\smallbreak\pagebreak[2]}
\correspDesc{Versand  durch Berta Zuckerkandl am [23. 6. 1911?] in Paris
\newline{}Erhalt  durch Arthur Schnitzler im Zeitraum [24. 6. 1911
                  – 28. 6. 1911?] in Wien}\toendnotes[C]{\smallbreak}
\Standort{CUL, Schnitzler, B 200.}
\physDesc{Brief, 2 Blätter, 7 Seiten, 1546 Zeichen (Absendeadresse »\textcolor{gray}{\textbf{84 RUE DE LONGCHAMP}}« auch auf dem zweiten Briefbogen)
\newline{}Handschrift: schwarze Tinte, lateinische Kurrent
\newline{}Schnitzler: mit Bleistift beschriftet: »\textsc{Zuckerkan\textcolor{gray}{dl}}« }\toendnotes[C]{\smallbreak}
\pstart
           {\pb}\uline{Freitag}.\hfill \textcolor{pink}{\textcolor{gray}{\textbf{84 RUE DE LONGCHAMP}}}\oindex{84, rue de Longchamp@\textbf{84, rue de Longchamp}, \emph{Wohngebäude}|pw}{}\ledrightnote{\textcolor{pink}{84, rue de Longchamp}}\pend
           
\pstart{}Verehrter Herr Doktor!\pend\vspace{0.5em}
\pstart
           \label{K_L03995-1v}\edtext{Gestern Donnerstag}{\lemma{\textnormal{\emph{Gestern Donnerstag}}}\Cendnote{\textnormal{\textcolor{blue}{Schnitzler} vermerkte für den 20. 6. 1911 im \emph{\textcolor{green}{Tagebuch}\pwindex{Schnitzler, Arthur 15. 5. 1862 Wien – 21. 10. 1931 ebd.@\textsc{Schnitzler, Arthur} (15. 5. 1862 Wien – 21. 10. 1931 ebd.), \emph{Schriftsteller, Mediziner}!Tagebuch@\strich\emph{Tagebuch}|pwk}} das Absenden des \emph{\textcolor{green}{Medardus}\pwindex{Schnitzler, Arthur 15. 5. 1862 Wien – 21. 10. 1931 ebd.@\textsc{Schnitzler, Arthur} (15. 5. 1862 Wien – 21. 10. 1931 ebd.), \emph{Schriftsteller, Mediziner}!junge Medardus. Dramatische Historie in einem Vorspiel und fünf Aufzügen@\strich\emph{Der junge Medardus. Dramatische Historie in einem Vorspiel und fünf Aufzügen}|pwk}}-\textcolor{green}{Szenariums}\pwindex{Schnitzler, Arthur 15. 5. 1862 Wien – 21. 10. 1931 ebd.@\textsc{Schnitzler, Arthur} (15. 5. 1862 Wien – 21. 10. 1931 ebd.), \emph{Schriftsteller, Mediziner}!junge Medardus. (Zur Übersetzung ins Französische)@\strich\emph{Der junge Medardus. (Zur Übersetzung ins Französische)}|pwkv} an \textcolor{blue}{Berta
                     Zuckerkandl}\pwindex{Zuckerkandl, Berta 13.\,4.\,1864 Wien – 16.\,10.\,1945 Paris@\textsc{Zuckerkandl, Berta} (13.\,4.\,1864 Wien – 16.\,10.\,1945 Paris), \emph{Schriftstellerin, Journalistin, Übersetzerin}|pwk} nach \textcolor{pink}{Paris}\oindex{Paris@\textbf{Paris}, \emph{Hauptstadt}|pwk}. Da zwei Tage
                  eine übliche Postlaufzeit für eine Sendung von \textcolor{pink}{Wien}\oindex{Wien@\textbf{Wien}, \emph{Verwaltungsgebiet}|pwk} nach \textcolor{pink}{Paris}\oindex{Paris@\textbf{Paris}, \emph{Hauptstadt}|pwk} darstellen, dürfte
                  der Zustelltermin der 22. 6. 1911 gewesen sein. Somit lässt sich der
                  vorliegende Brief auf den 23. 6. 1911 datieren.}}}\label{K_L03995-1} ist Ihr \textcolor{green}{Scenario}\pwindex{Schnitzler, Arthur 15. 5. 1862 Wien – 21. 10. 1931 ebd.@\textsc{Schnitzler, Arthur} (15. 5. 1862 Wien – 21. 10. 1931 ebd.), \emph{Schriftsteller, Mediziner}!junge Medardus. (Zur Übersetzung ins Französische)@\strich\emph{Der junge Medardus. (Zur Übersetzung ins Französische)}|pwv}{}\ledrightnote{{$\rightarrow$}\emph{\textcolor{green}{Der junge Medardus. (Zur Übersetzung ins Französische)}}} geko{\geminationm}en. Wofür ich herzlichst danke. Es ist ausgezeichnet
               gemacht – u ich beginne heute mit der Übersetzung.\pend
           
\pstart
           Für meine Unterre{\pb}dung mit den Direktoren
                  \textcolor{blue}{Herz}\pwindex{Hertz, Henri 17.\,6.\,1875 Nogent-sur-Seine – 11.\,10.\,1966 16. arrondissement [Paris]@\textsc{Hertz, Henri} (17.\,6.\,1875 Nogent-sur-Seine – 11.\,10.\,1966 16. arrondissement [Paris]), \emph{Schriftsteller, Journalist, Theaterdirektor}|pw}{}\ledrightnote{\textcolor{blue}{Henri Hertz}} und \textcolor{blue}{Coquelin}\pwindex{Coquelin, Jean 1.\,12.\,1865 Paris – 1.\,10.\,1944@\textsc{Coquelin, Jean} (1.\,12.\,1865 Paris – 1.\,10.\,1944), \emph{Schauspieler}|pw}{}\ledrightnote{\textcolor{blue}{Jean Coquelin}} kam es zu spät. Wieder alles Erwarten – da doch in
                  \textcolor{pink}{Paris}\oindex{Paris@\textbf{Paris}, \emph{Hauptstadt}|pw}{}\ledrightnote{\textcolor{pink}{Paris}} Alles so lange dauert – erhielt ich am
                     Mit{[}t{]}woch{ }Früh von der \textcolor{brown}{Porte St.
                  Martin}\orgindex{Théâtre de la Porte Saint-Martin@Théâtre de la Porte Saint-Martin|pwv}{}\ledrightnote{{$\rightarrow$}\emph{\textcolor{brown}{Théâtre de la Porte Saint-Martin}}} die \label{T_L03995-1v}\edtext{Verständigung}{\lemma{\textnormal{\emph{Verständigung}}}\Cendnote{\textnormal{Sie schreibt:
                  »Verständidung«.}}}\label{T_L03995-1} dass die \textcolor{blue}{Direktoren}\pwindex{Coquelin, Jean 1.\,12.\,1865 Paris – 1.\,10.\,1944@\textsc{Coquelin, Jean} (1.\,12.\,1865 Paris – 1.\,10.\,1944), \emph{Schauspieler}|pw}\pwindex{Hertz, Henri 17.\,6.\,1875 Nogent-sur-Seine – 11.\,10.\,1966 16. arrondissement [Paris]@\textsc{Hertz, Henri} (17.\,6.\,1875 Nogent-sur-Seine – 11.\,10.\,1966 16. arrondissement [Paris]), \emph{Schriftsteller, Journalist, Theaterdirektor}|pw}{}\ledrightnote{\textcolor{blue}{Jean Coquelin}{\newline}\textcolor{blue}{Henri Hertz}} mich bitten
                     Mit{[}t{]}woch{ }Nachmittag sie {\pb}zu besuchen.
               Ich hatte zum  Glück zur Vorsicht – einen langen Auszug aus dem \textcolor{green}{Medardus}\pwindex{Schnitzler, Arthur 15. 5. 1862 Wien – 21. 10. 1931 ebd.@\textsc{Schnitzler, Arthur} (15. 5. 1862 Wien – 21. 10. 1931 ebd.), \emph{Schriftsteller, Mediziner}!junge Medardus. Dramatische Historie in einem Vorspiel und fünf Aufzügen@\strich\emph{Der junge Medardus. Dramatische Historie in einem Vorspiel und fünf Aufzügen}|pw}{}\ledrightnote{\textcolor{green}{Der junge Medardus. Dramatische Historie in einem Vorspiel und fünf Aufzügen}} auf der Reise \textcolor{pink}{französisch}\oindex{Frankreich@\textbf{Frankreich}|pw}{}\ledrightnote{\textcolor{pink}{Frankreich}} geschrieben. Rasch diktirte ich dies in einem
               Schreib-Maschin-Bureau ab – und um 5 Uhr war ich bei \textcolor{blue}{Herz}\pwindex{Hertz, Henri 17.\,6.\,1875 Nogent-sur-Seine – 11.\,10.\,1966 16. arrondissement [Paris]@\textsc{Hertz, Henri} (17.\,6.\,1875 Nogent-sur-Seine – 11.\,10.\,1966 16. arrondissement [Paris]), \emph{Schriftsteller, Journalist, Theaterdirektor}|pw}{}\ledrightnote{\textcolor{blue}{Henri Hertz}}{ }{\kaufmannsund}{ }\textcolor{blue}{Coquelin}\pwindex{Coquelin, Jean 1.\,12.\,1865 Paris – 1.\,10.\,1944@\textsc{Coquelin, Jean} (1.\,12.\,1865 Paris – 1.\,10.\,1944), \emph{Schauspieler}|pw}{}\ledrightnote{\textcolor{blue}{Jean Coquelin}} – begleitet von D\textsuperscript{r}{ }\textcolor{blue}{{\pb}Frischauer}\pwindex{Frischauer, Berthold 9.\,9.\,1851 Brünn – 4.\,2.\,1924 Wien@\textsc{Frischauer, Berthold} (9.\,9.\,1851 Brünn – 4.\,2.\,1924 Wien), \emph{Journalist}|pw}{}\ledrightnote{\textcolor{blue}{Berthold Frischauer}}. \textcolor{blue}{Herz}\pwindex{Hertz, Henri 17.\,6.\,1875 Nogent-sur-Seine – 11.\,10.\,1966 16. arrondissement [Paris]@\textsc{Hertz, Henri} (17.\,6.\,1875 Nogent-sur-Seine – 11.\,10.\,1966 16. arrondissement [Paris]), \emph{Schriftsteller, Journalist, Theaterdirektor}|pw}{}\ledrightnote{\textcolor{blue}{Henri Hertz}} ist ein \label{K_L03995-2v}\edtext{\begin{otherlanguage}{french}Homme d’affaire\end{otherlanguage}}{\lemma{\textnormal{\emph{Homme d’affaire}}}\Cendnote{\textnormal{französisch: Geschäftsmann}}}\label{K_L03995-2} pur –
               der literarisch nicht viel weiss. \textcolor{blue}{Coquelin}\pwindex{Coquelin, Jean 1.\,12.\,1865 Paris – 1.\,10.\,1944@\textsc{Coquelin, Jean} (1.\,12.\,1865 Paris – 1.\,10.\,1944), \emph{Schauspieler}|pw}{}\ledrightnote{\textcolor{blue}{Jean Coquelin}}
               dagegen wusste Manches vom \textcolor{green}{Medardus}\pwindex{Schnitzler, Arthur 15. 5. 1862 Wien – 21. 10. 1931 ebd.@\textsc{Schnitzler, Arthur} (15. 5. 1862 Wien – 21. 10. 1931 ebd.), \emph{Schriftsteller, Mediziner}!junge Medardus. Dramatische Historie in einem Vorspiel und fünf Aufzügen@\strich\emph{Der junge Medardus. Dramatische Historie in einem Vorspiel und fünf Aufzügen}|pw}{}\ledrightnote{\textcolor{green}{Der junge Medardus. Dramatische Historie in einem Vorspiel und fünf Aufzügen}}. Ich wurde
               gebeten meine Inhalts-Angabe dort zu lassen. Aber ich gab mündlich eine Schilderung
               die wie mir D\textsuperscript{r} \textcolor{blue}{Frischauer}\pwindex{Frischauer, Berthold 9.\,9.\,1851 Brünn – 4.\,2.\,1924 Wien@\textsc{Frischauer, Berthold} (9.\,9.\,1851 Brünn – 4.\,2.\,1924 Wien), \emph{Journalist}|pw}{}\ledrightnote{\textcolor{blue}{Berthold Frischauer}} dann sagte – riesig in{[}ter{]}{\pb}ressirte. Besonders die \label{K_L03995-3v}\edtext{Pretendenten-Geschichte}{\lemma{\textnormal{\emph{Pretendenten-Geschichte}}}\Cendnote{\textnormal{Eine Nebenhandlung von \emph{\textcolor{green}{Der junge Medardus}\pwindex{Schnitzler, Arthur 15. 5. 1862 Wien – 21. 10. 1931 ebd.@\textsc{Schnitzler, Arthur} (15. 5. 1862 Wien – 21. 10. 1931 ebd.), \emph{Schriftsteller, Mediziner}!junge Medardus. Dramatische Historie in einem Vorspiel und fünf Aufzügen@\strich\emph{Der junge Medardus. Dramatische Historie in einem Vorspiel und fünf Aufzügen}|pwk}} umkreist den Anspruch auf die \textcolor{pink}{französische}\oindex{Frankreich@\textbf{Frankreich}|pwk} Krone, den in \textcolor{pink}{Wien}\oindex{Wien@\textbf{Wien}, \emph{Verwaltungsgebiet}|pwk} exilierte Adelsfamilie stellt.}}}\label{K_L03995-3} finden die \textcolor{blue}{Herrn}\pwindex{Coquelin, Jean 1.\,12.\,1865 Paris – 1.\,10.\,1944@\textsc{Coquelin, Jean} (1.\,12.\,1865 Paris – 1.\,10.\,1944), \emph{Schauspieler}|pwv}\pwindex{Hertz, Henri 17.\,6.\,1875 Nogent-sur-Seine – 11.\,10.\,1966 16. arrondissement [Paris]@\textsc{Hertz, Henri} (17.\,6.\,1875 Nogent-sur-Seine – 11.\,10.\,1966 16. arrondissement [Paris]), \emph{Schriftsteller, Journalist, Theaterdirektor}|pwv}{}\ledrightnote{{$\rightarrow$}\emph{\textcolor{blue}{Jean Coquelin}}{\newline}{$\rightarrow$}\emph{\textcolor{blue}{Henri Hertz}}} für \textcolor{pink}{Paris}\oindex{Paris@\textbf{Paris}, \emph{Hauptstadt}|pw}{}\ledrightnote{\textcolor{pink}{Paris}} höchst erwünscht. Meine Unterredung hatte
               wie es sich zeigte so rasch stattfinden müssen, weil \textcolor{blue}{Herz}\pwindex{Hertz, Henri 17.\,6.\,1875 Nogent-sur-Seine – 11.\,10.\,1966 16. arrondissement [Paris]@\textsc{Hertz, Henri} (17.\,6.\,1875 Nogent-sur-Seine – 11.\,10.\,1966 16. arrondissement [Paris]), \emph{Schriftsteller, Journalist, Theaterdirektor}|pw}{}\ledrightnote{\textcolor{blue}{Henri Hertz}}{ }{\kaufmannsund}{ }\textcolor{blue}{Coquelin}\pwindex{Coquelin, Jean 1.\,12.\,1865 Paris – 1.\,10.\,1944@\textsc{Coquelin, Jean} (1.\,12.\,1865 Paris – 1.\,10.\,1944), \emph{Schauspieler}|pw}{}\ledrightnote{\textcolor{blue}{Jean Coquelin}}{ }heute{ }{\pb}verreisen. Sie baten mich das von Ihnen
               zu sendende \textcolor{green}{Scenarium}\pwindex{Schnitzler, Arthur 15. 5. 1862 Wien – 21. 10. 1931 ebd.@\textsc{Schnitzler, Arthur} (15. 5. 1862 Wien – 21. 10. 1931 ebd.), \emph{Schriftsteller, Mediziner}!junge Medardus. Dramatische Historie in einem Vorspiel und fünf Aufzügen@\strich\emph{Der junge Medardus. Dramatische Historie in einem Vorspiel und fünf Aufzügen}|pwv}{}\ledrightnote{{$\rightarrow$}\emph{\textcolor{green}{Der junge Medardus. Dramatische Historie in einem Vorspiel und fünf Aufzügen}}}
               möglichst rasch zu übersetzen – {\kaufmannsund} an eine mir angegebene
               Adresse zu senden. Dann erst werden wir etwas wissen. Es ist möglich dass ich im
                  August von {\pb}der \textcolor{pink}{Schweiz}\oindex{Schweiz@\textbf{Schweiz}|pw}{}\ledrightnote{\textcolor{pink}{Schweiz}} wieder zu den \textcolor{blue}{Herrn}\pwindex{Coquelin, Jean 1.\,12.\,1865 Paris – 1.\,10.\,1944@\textsc{Coquelin, Jean} (1.\,12.\,1865 Paris – 1.\,10.\,1944), \emph{Schauspieler}|pw}\pwindex{Hertz, Henri 17.\,6.\,1875 Nogent-sur-Seine – 11.\,10.\,1966 16. arrondissement [Paris]@\textsc{Hertz, Henri} (17.\,6.\,1875 Nogent-sur-Seine – 11.\,10.\,1966 16. arrondissement [Paris]), \emph{Schriftsteller, Journalist, Theaterdirektor}|pw}{}\ledrightnote{\textcolor{blue}{Jean Coquelin}{\newline}\textcolor{blue}{Henri Hertz}} ko{\geminationm}en muss.
               Jedenfalls war die Anknüpfung günstig. Mehr kann ich nicht sagen. Einige Détails noch
                  morgen oder übermorgen – da ich heute rasch
               schliessen muss.\pend
           
\pstart
           Ihnen u Ihrer \textcolor{blue}{Frau}\pwindex{Schnitzler, Olga 17.\,1.\,1882 Wien – 13.\,1.\,1970 Lugano@\textsc{Schnitzler, Olga} (17.\,1.\,1882 Wien – 13.\,1.\,1970 Lugano), \emph{Schauspielerin, Sängerin}|pwv}{}\ledrightnote{{$\rightarrow$}\emph{\textcolor{blue}{Olga Schnitzler}}}
               herzliche Grüsse.{\\[\baselineskip]} Ihre \spacefill\mbox{Berta Zuckerkandl}\pend
           \leftskip=0em{}\selectlanguage{ngerman}\endnumbering\briefempfaengerindex{Schnitzler, Arthur@\textsc{Schnitzler, Arthur}!zzzZuckerkandl, Berta@\emph{von Berta Zuckerkandl}!1911-06-231@{{[}23. 6. 1911?{]}}|)be}\mylabel{L03995h}  \newcommand{\dateiname}{L03995}\newcommand{\titel}{Berta Zuckerkandl an Arthur Schnitzler, [23. 6. 1911?]}\newcommand{\editorInnen}{Herausgegeben von Jahnke, SelmaMüller, Martin Anton}\input{../tex-inputs/latex-korrekturansicht-abspann}
